\chapter{Higher direct images of quasi-coherent sheaves in positive degree}
\label{chap:Cohomology_HigherDirectImage}
% archon:covers AlgebraicJacobian/Cohomology/HigherDirectImage.lean

\section{Setup and motivation}

This chapter treats the higher derived direct images \(R^i f_*\mathcal{F}\) for
\(i \geq 1\), complementing the \(i = 0\) (direct-image) case of
Chapter~\ref{chap:Cohomology_FlatBaseChange}. These higher direct images are the
single deepest foundation of the cohomology-and-base-change package: the
Cohen--Macaulay regularity bound, semicontinuity of fibre cohomology, and the
\(m\)-regularity estimate of FGA all rest on the basic finiteness and base-change
properties of \(R^i f_*\) recorded here.

We use the classical cohomology-sheaf formulation (not the derived category) and a
{\v C}ech-complex approach, parallel to the \(i = 0\) treatment of
Chapter~\ref{chap:Cohomology_FlatBaseChange}.

Throughout, \(f : X \to S\) is a morphism of schemes that is quasi-compact and
quasi-separated, and \(\mathcal{F}\) is a quasi-coherent \(\mathcal{O}_X\)-module.
A base change of \(f\) along a morphism \(g : S' \to S\) is recorded by the
cartesian square obtained from the fibre product \(X' = X \times_S S'\):
\[
  \begin{array}{ccc}
    X' & \xrightarrow{\;g'\;} & X \\[2pt]
    {\scriptstyle f'}\big\downarrow & & \big\downarrow{\scriptstyle f} \\[2pt]
    S' & \xrightarrow{\;g\;} & S
  \end{array}
\]
Here \(f' : X' \to S'\) is the base change of \(f\), \(g' : X' \to X\) is the first
projection, and \(\mathcal{F}' = (g')^*\mathcal{F}\) is the pullback of
\(\mathcal{F}\) to \(X'\).

\section{Definition of the higher direct images}

\begin{definition}[Higher direct image]\leanok
  \label{def:higher_direct_image}
  \lean{AlgebraicGeometry.higherDirectImage}
  \dcref{custom}
  Let \(f : X \to S\) be a morphism of schemes and let \(\mathcal{F}\) be a
  quasi-coherent \(\mathcal{O}_X\)-module. For \(i \geq 0\) the \emph{\(i\)-th
  higher direct image} \(R^i f_*\mathcal{F}\) is the \(i\)-th right derived functor
  of the pushforward functor \(f_*\) applied to \(\mathcal{F}\). Concretely,
  \(R^i f_*\mathcal{F}\) is the \(\mathcal{O}_S\)-module obtained as the
  sheafification of the presheaf
  \[
    V \;\longmapsto\; H^i\bigl(f^{-1}(V),\, \mathcal{F}|_{f^{-1}(V)}\bigr),
    \qquad V \subset S \text{ open},
  \]
  where \(H^i(-, -)\) denotes sheaf cohomology. For \(i = 0\) this recovers the
  ordinary pushforward \(R^0 f_*\mathcal{F} = f_*\mathcal{F}\).
\end{definition}

\section{Quasi-coherence of the higher direct images}

\begin{lemma}[Quasi-coherence of higher direct images]
  \label{lem:higher_direct_image_quasi_coherent}
  \uses{def:higher_direct_image, lem:higher_direct_image_affine_vanishing}
  \dcref{custom}
  Let \(f : X \to S\) be a quasi-compact and quasi-separated morphism of schemes.
  Then for every quasi-coherent \(\mathcal{O}_X\)-module \(\mathcal{F}\) and every
  \(p \geq 0\), the higher direct image \(R^p f_*\mathcal{F}\) is a quasi-coherent
  \(\mathcal{O}_S\)-module.
\end{lemma}

\begin{proof}
  \uses{def:higher_direct_image, lem:higher_direct_image_affine_vanishing}
  Quasi-coherence is local on \(S\), and forming higher direct images commutes
  with restriction to open subschemes, so we may assume \(S\) is affine; then
  \(X\) is quasi-compact and quasi-separated. We argue by the induction principle
  for quasi-compact opens of a quasi-compact quasi-separated scheme: for a
  quasi-compact open \(U \subset X\) with \(a = f|_U\), let \(P(U)\) be the
  property that \(R^p a_*\mathcal{G}\) is quasi-coherent on \(S\) for every
  quasi-coherent module \(\mathcal{G}\) on \(U\) and every \(p \geq 0\); the
  desired statement is \(P(X)\).

  \emph{Affine base.} If \(U\) is affine then \(a\) is an affine morphism, so
  \(R^p a_*\mathcal{G} = 0\) for \(p \geq 1\) by the relative affine vanishing of
  Lemma~\ref{lem:higher_direct_image_affine_vanishing}, while
  \(R^0 a_*\mathcal{G} = a_*\mathcal{G}\) is quasi-coherent because the pushforward
  of a quasi-coherent sheaf along a quasi-compact quasi-separated morphism is
  quasi-coherent. Hence \(P(U)\) holds.

  \emph{Inductive step.} Let \(U \subset X\) be quasi-compact open, \(V \subset X\)
  affine open, and suppose \(P(U)\), \(P(V)\), \(P(U \cap V)\) hold. Writing
  \(a = f|_U\), \(b = f|_V\), \(c = f|_{U \cap V}\), \(h = f|_{U \cup V}\), the
  relative Mayer--Vietoris sequence
  \[
    0 \to h_*\mathcal{G} \to a_*(\mathcal{G}|_U) \oplus b_*(\mathcal{G}|_V)
      \to c_*(\mathcal{G}|_{U \cap V}) \to R^1 h_*\mathcal{G} \to \cdots
  \]
  exhibits each \(R^p h_*\mathcal{G}\) as an extension built from quasi-coherent
  sheaves \(R^p a_*(\mathcal{G}|_U)\), \(R^p b_*(\mathcal{G}|_V)\),
  \(R^p c_*(\mathcal{G}|_{U \cap V})\): it sits between the cokernel of a map of
  quasi-coherent sheaves and the kernel of a map of quasi-coherent sheaves, hence
  is quasi-coherent. Thus \(P(U \cup V)\) holds. By the induction principle
  \(P(X)\) holds, proving the lemma.
\end{proof}

\section{Vanishing of higher direct images along an affine morphism}

\begin{lemma}[Relative affine vanishing]
  \label{lem:higher_direct_image_affine_vanishing}
  \uses{def:higher_direct_image}
  \dcref{custom}
  Let \(f : X \to S\) be a morphism of schemes and \(\mathcal{F}\) a
  quasi-coherent \(\mathcal{O}_X\)-module. If \(f\) is affine, then
  \(R^i f_*\mathcal{F} = 0\) for all \(i \geq 1\).
\end{lemma}

\begin{proof}
  \uses{def:higher_direct_image}
  By Definition~\ref{def:higher_direct_image}, \(R^i f_*\mathcal{F}\) is the
  sheafification of the presheaf \(V \mapsto H^i(f^{-1}(V),
  \mathcal{F}|_{f^{-1}(V)})\). Since \(f\) is affine, the preimage \(f^{-1}(V)\) of
  any affine open \(V \subset S\) is affine. The higher cohomology of a
  quasi-coherent sheaf on an affine scheme vanishes, so
  \(H^i(f^{-1}(V), \mathcal{F}|_{f^{-1}(V)}) = 0\) for all \(i \geq 1\) and all
  affine \(V\). As the affine opens form a basis of \(S\), the sheafification of a
  presheaf that vanishes on a basis is zero; hence \(R^i f_*\mathcal{F} = 0\) for
  all \(i \geq 1\).
\end{proof}

\section{Flat base change for the higher direct images}

\begin{theorem}[Flat base change, $i \geq 1$ case]
  \label{thm:flat_base_change_higher}
  \uses{def:higher_direct_image, lem:higher_direct_image_quasi_coherent,
    def:pushforward_base_change_map, lem:cech_computes_cohomology,
    lem:cech_flat_base_change}
  \dcref{custom}
  Consider the cartesian square
  \[
    \begin{array}{ccc}
      X' & \xrightarrow{\;g'\;} & X \\[2pt]
      {\scriptstyle f'}\big\downarrow & & \big\downarrow{\scriptstyle f} \\[2pt]
      S' & \xrightarrow{\;g\;} & S
    \end{array}
  \]
  with \(\mathcal{F}\) a quasi-coherent \(\mathcal{O}_X\)-module and
  \(\mathcal{F}' = (g')^*\mathcal{F}\). Assume that \(g\) is flat and that \(f\) is
  quasi-compact and quasi-separated. Then for every \(i \geq 1\) the canonical
  base-change map (the higher-degree analogue of the pushforward base-change map
  of Definition~\ref{def:pushforward_base_change_map}) is an isomorphism
  \[
    g^*\bigl(R^i f_*\mathcal{F}\bigr) \;\xrightarrow{\;\sim\;}\;
    R^i f'_*\bigl((g')^*\mathcal{F}\bigr) = R^i f'_*\mathcal{F}' .
  \]
  Equivalently, in the affine situation \(S = \operatorname{Spec}(A)\),
  \(S' = \operatorname{Spec}(B)\) with \(A \to B\) flat, the comparison map of
  \(B\)-modules
  \[
    H^i(X, \mathcal{F}) \otimes_A B \;\xrightarrow{\;\sim\;}\;
    H^i(X_B, \mathcal{F}_B)
  \]
  is an isomorphism, where \(X_B = X \times_{\operatorname{Spec}(A)}
  \operatorname{Spec}(B)\) and \(\mathcal{F}_B\) is the pullback of
  \(\mathcal{F}\).
\end{theorem}

\begin{proof}
  \uses{def:higher_direct_image, lem:higher_direct_image_quasi_coherent,
    def:pushforward_base_change_map}
  \emph{Reduction to the affine target.} The assertion is local on \(S'\), so we
  may assume \(S = \operatorname{Spec}(A)\) and \(S' = \operatorname{Spec}(B)\) with
  \(A \to B\) flat. By quasi-coherence of the higher direct images
  (Lemma~\ref{lem:higher_direct_image_quasi_coherent}), \(R^i f_*\mathcal{F}\) is
  the quasi-coherent \(\mathcal{O}_S\)-module associated to the \(A\)-module
  \(H^0(S, R^i f_*\mathcal{F}) = H^i(X, \mathcal{F})\), and likewise
  \(R^i f'_*\mathcal{F}'\) is associated to the \(B\)-module
  \(H^i(X', \mathcal{F}')\). Since pullback along \(g\) corresponds to
  \(- \otimes_A B\) on modules, the base-change map is identified with the
  comparison map of \(B\)-modules
  \[
    H^i(X, \mathcal{F}) \otimes_A B \longrightarrow H^i(X_B, \mathcal{F}_B),
    \qquad X_B = X \times_{\operatorname{Spec}(A)} \operatorname{Spec}(B),
  \]
  and it suffices to prove that this map is an isomorphism.

  \medskip
  \emph{The separated case via {\v C}ech complexes.} Suppose first that \(X\) is
  separated over \(A\). Choose a finite affine open cover
  \(\mathcal{U} : X = U_1 \cup \dots \cup U_t\); by separatedness every finite
  intersection \(U_{i_0 \dots i_p}\) is again affine, so {\v C}ech cohomology
  computes sheaf cohomology:
  \(\check{H}^p(\mathcal{U}, \mathcal{F}) = H^p(X, \mathcal{F})\) for all \(p\).
  Base changing the cover yields an affine open cover
  \(\mathcal{U}_B : X_B = (U_1)_B \cup \dots \cup (U_t)_B\) of the again-separated
  \(X_B\), with the same property
  \(\check{H}^p(\mathcal{U}_B, \mathcal{F}_B) = H^p(X_B, \mathcal{F}_B)\). The
  affine case identifies each term of the base-changed {\v C}ech complex with the
  original term tensored over \(A\) with \(B\); that is, there is an isomorphism of
  cochain complexes
  \[
    \check{\mathcal{C}}^\bullet(\mathcal{U}_B, \mathcal{F}_B) \;\cong\;
    \check{\mathcal{C}}^\bullet(\mathcal{U}, \mathcal{F}) \otimes_A B .
  \]
  Because \(A \to B\) is flat, the functor \(- \otimes_A B\) is exact and hence
  commutes with the formation of cohomology of a complex; taking \(H^i\) in any
  degree \(i \geq 1\) gives the desired isomorphism
  \(H^i(X, \mathcal{F}) \otimes_A B \cong H^i(X_B, \mathcal{F}_B)\).

  \medskip
  \emph{The quasi-separated case.} In general \(X\) need only be quasi-separated.
  Choose a finite affine open cover \(\mathcal{U} : X = U_1 \cup \dots \cup U_t\).
  Now the intersections \(U_{i_0 \dots i_p}\) are quasi-compact and separated rather
  than affine, but they fall under the separated case already treated, so flat base
  change holds for each of them. Comparing the two {\v C}ech-to-cohomology spectral
  sequences with \(E_2\)-pages
  \(\check{H}^p(\mathcal{U}, \underline{H}^q(\mathcal{F}))\) and
  \(\check{H}^p(\mathcal{U}_B, \underline{H}^q(\mathcal{F}_B))\), converging to
  \(H^{p+q}(X, \mathcal{F})\) and \(H^{p+q}(X_B, \mathcal{F}_B)\) respectively, the
  separated case provides the term-by-term identification
  \(\check{H}^p(\mathcal{U}_B, \underline{H}^q(\mathcal{F}_B)) =
  \check{H}^p(\mathcal{U}, \underline{H}^q(\mathcal{F})) \otimes_A B\) on the
  \(E_2\)-pages; flatness of \(A \to B\) propagates this isomorphism through the
  spectral sequences to the abutment, giving
  \(H^i(X, \mathcal{F}) \otimes_A B \cong H^i(X_B, \mathcal{F}_B)\) in every degree
  \(i \geq 1\). Translating back to sheaves on \(S'\) yields the asserted
  isomorphism \(g^*(R^i f_*\mathcal{F}) \xrightarrow{\sim} R^i f'_*\mathcal{F}'\).
\end{proof}
